\section{Reflection on Readiness}
\label{sec:readiness}

For a runner in Charlottesville who looks at a route before running it,
RunLine is ready to use. The evidence for that is strong and repeatable. Every
benchmark request in both modes returned routes, distance fit met its
criterion, interruptions fell by more than 70\% against a distance-only
planner, the surface preference moved routes the right way in 97 to 98\% of
cases, every GPX file we generated validated, and warm requests answer in
under a second. The precomputed routes served in production are on average
better than live planning: all were within the distance band and fewer
repeated themselves. The six additional metro areas passed the same smoke
checks, which shows the pipeline generalizes beyond the city it was tuned on.

What remains uncertain is mostly at the edges of that use. Until 18 November the engine was ahead of the interface: a runner saw a good
route but not why it was chosen, which is the part of the design the proposal
said distinguished RunLine. That display is now in place and verified, but no
runner outside the team has used it yet. The elevation preference now works
for live requests too, but its fix is a single tuned constant measured on
eight starts, and the hilly setting still relies on ranking alone. The
production backend now runs the code we measured, but the regional server for
the metro areas was not part of that release. Its first request in a new city
is now fast when the page has pre-loaded the map, but still takes several
seconds when it has not. We also have no evidence about
people other than ourselves, and none about the GPS run mode outdoors.
Accessibility has an automated score of 96 but no keyboard or screen-reader
test.

Before wider use, three things must happen, in order. We must complete the remaining browser checklists and the route review;
name the map's start pin and do a keyboard and screen-reader pass; and fix
whatever those turn up, starting with the silent hilly result. Before anyone outside the team is invited to use it, we
would also want a short observed session with a handful of runners. The
proposal's premise was that visible tradeoffs help a runner choose, and
nothing in this report tests it. The metro
areas should not be promoted beyond a preview until Washington, New York and
Los Angeles are rebuilt with the crosswalk signals and the first request in a
new city is reliably fast.

Validation also changed how we work, through three lessons. The first is that a metric needs a check
that does not share its assumptions. Building the independent crossing
labeller found errors in both directions: the labeller was wrong in eight of
the cases we first blamed on the planner, and the planner was missing a whole
class of crossings. Without the audit, the US3 result would rest on a count
nobody had checked. The second is that tests prove what they assert and
nothing else. Our interface tests confirmed that every displayed number was
correct, and so stayed green when the numbers that mattered most for US4 were
removed from the screen. Measuring against the proposal's criteria, rather
than against our tests, is what exposed it, and it is why our browser tests now
check that required fields are present and why tests run on every pull request. The third
is that the deployed system is its own thing to validate. The backend was
behind the code we benchmarked, and a caching behaviour broke our own
rehearsal scripts while leaving browsers unaffected. Neither shows up on a
laptop. The validation plan for the final sprint therefore reruns the
acceptance and deployment scripts against production after every release,
not only the test suite before it.
